\documentclass[10pt,a4paper]{amsart}
\usepackage[margin=19mm]{geometry}
\usepackage{amsmath,amssymb,mathtools}
\usepackage[scr=rsfs,cal=euler]{mathalfa}
\usepackage{xfrac}
\usepackage[unicode,hidelinks,bookmarks=false]{hyperref}
\newcommand{\ie}{i.e.\ }
\newcommand{\cf}{cf.\ }
\newcommand{\resp}{respectively\ }
\newcommand{\Subsection}{\subsection}
\providecommand{\nobreakdash}{\nobreak\hbox{-}}
\setlength{\emergencystretch}{2em}
\multlinegap0pt
\usepackage{bm}
\usepackage{bbm}
\usepackage{dsfont}
\usepackage{xspace}
\usepackage{cancel}
\usepackage{mathtools}
\usepackage{amsthm}
\makeatletter
\@ifundefined{theorem}{\newtheorem{theorem}{Theorem}}{}
\@ifundefined{lemma}{\newtheorem{lemma}[theorem]{Lemma}}{}
\makeatother
\title[A Generalized Hacker Dynamics Model with Nonlinear Incidence]{A Generalized Hacker Dynamics Model with Nonlinear Incidence: Analysis and Positivity-Preserving Numerical Simulation}
\author{Manh Tuan Hoang, Hoai Thu Pham, Matthias Ehrhardt}
\date{Source: arXiv:2606.14191v1 (CC BY 4.0)}
\begin{document}
\maketitle

\noindent\emph{Source and licence.} A Generalized Hacker Dynamics Model with Nonlinear Incidence: Analysis and Positivity-Preserving Numerical Simulation. Version arXiv:2606.14191v1 (CC BY 4.0), distributed under the Creative Commons Attribution 4.0 International licence, as recorded in the arXiv licence metadata for this exact version and shown on the abstract page https://arxiv.org/abs/2606.14191v1. Full rights notice: licenses/arxiv-2606.14191-CC-BY-4.0.txt.\par\smallskip

\noindent\textit{Editorial scope.} Counted unit: the Lemma whose source label is \texttt{Lemma1} in the article (source0.tex lines 431--454, source label \texttt{Lemma1}), delivered together with the article's own complete original proof of it (source0.tex lines 456--504, one \texttt{proof} environment of 49 lines). No mathematics is written, repaired or re-derived: statement and proof are the article's own text line for line. Required context retained uncounted: eq:1G (lines 247--253). Every definition, display and label of those lines is retained unchanged. Omitted from this package and not redistributed: 1--246, 254--430, 455--455, 505--1592. Nothing inside a retained excerpt is omitted or altered: every retained body is delivered exactly as the article's lines are written, so no omission receipt of any kind accompanies this package. Citations: the retained excerpts carry no citation command at all, so no bibliography entry of the article is transcribed and no reference is redistributed. Reference adaptation: none was needed and none was made; every reference inside the delivered unit resolves inside it, so no unresolved reference or citation is produced. Printed slips kept literal: no printed word, symbol, hypothesis or number of the counted statement or of its proof was altered. Layout and class: the article's own class is \texttt{article} with packages including inputenc, fontenc, authblk, amsmath, amsfonts, amssymb, ulem, xcolor, amsthm, graphicx, marginnote, subcaption, hyperref, geometry; the wrapper is the lane's pinned \texttt{amsart} editorial wrapper at 10pt on A4, which restates the theorem-like environments the retained text needs and adds the article's own macro definitions that the retained text needs (none), with extra wrapper packages (bm, bbm, dsfont, xspace, cancel, mathtools). The delivered numbering is the unit's own. Assets: no retained span contains a figure environment, a graphics-inclusion command, a PDF-page inclusion, a table or a TikZ picture; no image file is copied into this package, the package declares no asset dependency and no image file is redistributed with it. Licence: version arXiv:2606.14191v1 of this article is distributed under the Creative Commons Attribution 4.0 International licence, as recorded in the arXiv licence metadata for this exact version and shown on the abstract page https://arxiv.org/abs/2606.14191v1. A full rights notice is delivered at licenses/arxiv-2606.14191-CC-BY-4.0.txt. Characters: every retained excerpt is pure ASCII, so the delivered engine, XeLaTeX with its default Latin Modern text font, reports no missing character.
\begin{equation}\label{eq:1G}
\begin{split}
 \dfrac{dC}{dt} &= \Lambda - f(C, H_n) -\alpha C + \gamma H_n -\mu C, \\
 \dfrac{dH_n}{dt} &= f(C, H_n) - \gamma H_n - \theta H_n - \mu H_n,\\
 \dfrac{d H_b}{dt} &= \alpha C + \theta H_n - \mu H_b.
\end{split}
\end{equation}

\begin{lemma}\label{Lemma1}
There are two types of equilibrium points of \eqref{eq:1G}:
\begin{itemize} 
\item[(iii)] A hacker-present equilibrium (HPE) point $E^* = \big(C^*,\,H_n^*,\,H_b^*\big)$ exists if and only if 
\begin{equation}\label{eq:14}
     \dfrac{\partial f\bigg(\dfrac{\Lambda}{\alpha + \mu},\,0\bigg)}{\partial H_n} - (\gamma + \theta + \mu) > 0.
\end{equation}
When this is the case, $E^*$ is defined by
\begin{equation}\label{eq:11}
\begin{split}
  H_b^* &= \dfrac{\alpha C^* + \theta H_n^*}{\mu},\\
  C^* &= \dfrac{\Lambda - (\theta + \mu)H_n^*}{\alpha + \mu}, 
\end{split}
\end{equation}
where $H_n^*$ is a unique positive solution, belonging to $\big(0,\,\frac{\Lambda}{\theta + \mu}\big)$, of the equation
\begin{equation}\label{eq:12}
    W(H_n) := \dfrac{1}{H_n}f\bigg(\dfrac{\Lambda - (\theta + \mu)H_n}{\alpha + \mu}, H_n\bigg) - (\gamma + \theta + \mu) = 0.
\end{equation}
\item[(ii)] A hacker-free equilibrium (HFE) point $E^0$ always exists, where
\begin{equation}\label{eq:15}
E^0 = \big(C^0,\,H_n^0,\,H_b^0\big) = \bigg(\dfrac{\Lambda}{\alpha + \mu},\,0,\,\dfrac{\alpha\Lambda}{\mu(\alpha + \mu)}\bigg).
\end{equation}
\end{itemize}
\end{lemma}

\begin{proof}
Any equilibrium point is a solution to the system
\begin{equation}\label{eq:10}
\begin{split}
0 &= \Lambda - f(C, H_n) -\alpha C + \gamma H_n -\mu C, \\
0 &= f(C, H_n) - \gamma H_n - \theta H_n - \mu H_n,\\
0 &= \alpha C + \theta H_n - \mu H_b.
\end{split}
\end{equation}
From the third equation of \eqref{eq:10}, we obtain the first formula of \eqref{eq:11}. 
On the other hand, adding side-by-side the first and second equations of \eqref{eq:10} gives
\begin{equation*}
    \Lambda - (\alpha + \mu)C - (\theta + \mu)H_n = 0,
\end{equation*}
which leads to the second formula of \eqref{eq:11}.

By substituting the formula of $C^*$ in the second formula of \eqref{eq:11} into the the second equation of \eqref{eq:10}, we obtain an equation for $H_n$
\begin{equation*}
    f\bigg(\dfrac{\Lambda - (\theta + \mu)H_n}{\alpha + \mu}, H_n\bigg) - \gamma H_n - \theta H_n - \mu H_n,
\end{equation*}
which is equivalent to $H_nW(H_n) = 0$, where $W$ is defined as in \eqref{eq:12}

Thus, any HPE point must satisfy \eqref{eq:12} and then, $H_n^* \in \bigl(0,\,\frac{\Lambda}{\theta + \mu}\bigr)$. 
Therefore, we consider the function $W$ on $\big(0,\,\frac{\Lambda}{\theta + \mu}\big)$.
On this interval, we have
\begin{equation}\label{eq:13}
\begin{split}
W\bigg(\dfrac{\Lambda}{\theta + \mu}\bigg) &= f\bigg(0,\, \dfrac{\Lambda}{\theta + \mu}\bigg) - (\gamma + \theta + \mu) = -(\gamma + \theta + \mu) < 0,\\
\lim_{H_n \to 0}\bigl(W(H_n)\bigr) &= \dfrac{\partial f\bigg(\dfrac{\Lambda}{\alpha + \mu},\,0\bigg)}{\partial H_n} - (\gamma + \theta + \mu) > 0.
\end{split}
\end{equation}
Note that the first estimate of \eqref{eq:13} is derived from the property $\textbf{(P1)}$, whereas the second one is derived from \eqref{eq:14}.

Furthermore, the first derivative of $W(H_n)$ with $H_n>0$ is given by
\begin{equation*}
\begin{split}
   W'(H_n) &= \dfrac{1}{H_n^2}\bigg[H_n\bigg(-\dfrac{\theta + \mu}{\alpha + \mu}\bigg)\dfrac{\partial f}{\partial C}\bigg(\dfrac{\Lambda - (\theta + \mu)H_n}{\alpha + \mu},\,H_n\bigg)\\
    &\qquad+ H_n\dfrac{\partial f}{\partial H_n}\bigg(\dfrac{\Lambda - (\theta + \mu)H_n}{\alpha + \mu},\,H_n\bigg) - f\bigg(\dfrac{\Lambda - (\theta + \mu)H_n}{\alpha + \mu},\,H_n\bigg)\bigg].
\end{split}
\end{equation*}
From $\textbf{(P2)}$ and $\textbf{(P4)}$, we deduce that $W'(H_n) < 0$ for $H_n \in \big(0,\,\frac{\Lambda}{\theta + \mu}\big)$. 
Combining this with \eqref{eq:13}, we conclude that the equation $W(H_n) = 0$  has a unique solution belonging to $\big(0,\,\frac{\Lambda}{\theta + \mu}\big)$. 
Therefore, $E^*$ is the only HPE point of \eqref{eq:1G}.

From the equation $H_nW(H_n) = 0$, we obtain the trivial solution $H_n^0 = 0$. 
Using \eqref{eq:12} yields $H_b^0$ and $C^0$, which are defined as in  \eqref{eq:15}. 
Thus, $E^0$ is an HFE point, which always exists for all parameter values. 
The proof is complete.
\end{proof}

\end{document}
